Un petit satellite sphérique est en orbite près de la Terre, au soleil (intensité solaire de \qty{1361}{W/m^2}). Son revêtement a une absorptivité de \num{\v{ab}} pour le rayonnement solaire et une émissivité de \num{\v{em}} dans l'infrarouge. On néglige le rayonnement reçu de la Terre et la puissance produite par l'électronique.

\begin{enumerate}[label=\alph*)]
    \item Écrivez le bilan de puissance du satellite et montrez que sa température d'équilibre ne dépend pas de son rayon.
    \item Déterminez la température d'équilibre du satellite.
    \item Quelle serait la température si le satellite était un corps noir ? Comment les ingénieurs utilisent-ils le choix du revêtement ?
\end{enumerate}
